\appendix
\section{Probability Identities and Approximation Limits}\label{app:proofs}
\subsection{Protection benefit and equal-budget swaps}
Let $G$ exclude an unprotected bit $b$. Independence gives
\[
J_p(\protset)=(1-p)\E R(G)+p\E R(G\cup\{b\}),\qquad
J_p(\protset\cup\{b\})=\E R(G).
\]
Subtracting proves Eq.~\eqref{eq:protection-benefit}. For the swap in Eq.~\eqref{eq:swap}, the old mask leaves $b$ exposed and the new mask leaves $a$ exposed. Both share the same Bernoulli background and no-fault term; subtracting yields the stated result. These are elementary conditioning identities. They do not require additive loss, differentiability or a neural network. They also do not identify naturally absent faults with interventions under arbitrary correlated or fixed-count laws.

For additive risk $R(A)=R(\varnothing)+\sum_{b\in A}I_b$, linearity gives $J_p(\protset)=R(\varnothing)+p\sum_{b\notin\protset}I_b$. Equal-cost, equal-probability exact-budget optimization therefore selects the largest $I_b$. The nonadditive counterexample in Sec.~\ref{sec:allocation} is nonnegative on all eight states and changes the optimal budget-one mask at $p=1/15$. It proves failure of a general guarantee, not failure for every network.

\subsection{Exact candidate interactions and a useful caution about bounds}
For a fixed outside field, the Boolean response has the unique multilinear expansion~\cite{odonnell2014analysis}
\[
h_G(A)=\sum_{T\subseteq A}\mu_G(T),\qquad
\mu_G(T)=\sum_{U\subseteq T}(-1)^{|T|-|U|}h_G(U).
\]
Thus $J_G(Q)=\sum_{T\subseteq C\setminus Q}p^{|T|}\mu_G(T)$, and the order-one truncation is precisely the outside-conditioned additive part of Eq.~\eqref{eq:decomposition}. Signed coefficients can cancel. A sum of absolute coefficient magnitudes is a bound or diagnostic, not the signed protection benefit. Boolean multilinear expansions and their expectation identities are established mathematics, not new foundations introduced by NetSuRF.

Suppose $h_G(A)\in[a,a+M]$ and $n=|C\setminus Q|$. The order-one prediction for a state with $s\ge2$ candidate faults is $\sum_{b\in A}h_G(\{b\})-(s-1)h_G(\varnothing)$. Its deviation from $h_G(A)$ is at most $sM$, while it is exact at $s=0,1$. For $S\sim\mathrm{Binomial}(n,p)$,
\begin{equation}\label{eq:bound}
\left|J_G(Q)-h_G(\varnothing)-p\sum_{b\in C\setminus Q}d_b^G\right|
\le M\E[S\mathbf1_{S\ge2}]
=Mnp\left[1-(1-p)^{n-1}\right].
\end{equation}
This gives $O(n(n-1)p^2)$ behavior at fixed small $n$ and low $p$. AUC risk has $M=1$; a valid bounded-loss range is needed for other tasks. The bound may be far too loose to certify an allocation because observed decision gaps can be much smaller. It bounds residual \emph{within-candidate} interactions after conditioning on the outside field; it says nothing comparable about the error of clean isolated scores caused by a large faulty background. Uniform risk approximation error $\varepsilon$ would imply at most $2\varepsilon$ excess risk for its minimizer over the same feasible set, but the present empirical studies do not establish a useful whole-network $\varepsilon$.

\section{Additional Evidence and Reproducibility Boundaries}\label{app:evidence}
The companion figure supplement supplies all four measured interior budgets for all six models. Machine-readable condition means, paired contrasts, cohort identities, source manifests, score-construction receipts and figure scripts are distributed separately. Each table or panel is traced to a finalized scientific payload; checksums and development identifiers are kept out of the reader-facing plots.

The earlier conditional attribution study includes a 12-pool KWS study at BER 0.01, a separate 12-pool KWS sweep at 0.001, 0.003 and 0.01, four ResNet-8 pools at 0.01, and four ToyCar pools at 0.001. Pool identities are reused across BER where specified and must not be counted as independent samples of architecture. Table~\ref{tab:attribution} compares strict mask-order reversals against the full response-table selection objective. Counts are dependent pairwise order comparisons, not independent statistical trials. All six saved across-pool held-out selected-minus-isolated intervals include zero.

\begin{table}[h]
\caption{Selection-order reversals relative to the exact response table. Conditioning individual candidate effects on outside faults removes most recorded reversals in these bounded designs. These counts do not establish held-out superiority.}\label{tab:attribution}
\centering\small
\begin{tabular}{llrr}\toprule
Model & BER & Clean additive & Outside-conditioned additive\\\midrule
KWS, initial pools & 0.01 & 396 & 8\\
KWS, multi-BER pools & 0.001 & 328 & 0\\
KWS, multi-BER pools & 0.003 & 305 & 0\\
KWS, multi-BER pools & 0.01 & 376 & 4\\
ResNet-8 & 0.01 & 201 & 1\\
ToyCar & 0.001 & 196 & 0\\\bottomrule
\end{tabular}
\Description{Full selection-order attribution for six historical model and BER groups, preserving negative and unresolved held-out findings.}
\end{table}

Additional bounded scaling studies on KWS and ResNet-8 examined candidate sizes 6, 12, 32 and 64 at three expected-fault levels, with four pools per condition. Exact selection was feasible only at the smaller sizes. At sizes 6 and 12, the outside-conditioned additive choice had zero recorded selection regret against the restricted exact choice in all 24 conditions per model. This is finite-table agreement, not a theorem about larger pools. All 54 simultaneous equal-pool intervals per model included zero; several differences were identically zero, which does not establish population equivalence. These measurements do not support a scalable optimizer-superiority claim. They help identify where expensive exact enumeration currently adds little observed decision value.

\subsection{Historical estimators and precision-specific studies}
Separate focused follow-ups strengthen two partial-tier comparisons without being pooled into the main grid. On ResNet-8 at $p=0.01$ and 6.25\% coverage, 320 paired fields give FKeras minus MSB cross-entropy $+2.241$ nats/image, interval $[1.890,2.593]$, in the same direction as the main study. On ToyCar at $p=0.01$ and 18.75\% coverage, 384 paired fields give FKeras minus absolute-gradient AUC risk $+0.1106$, interval $[0.0920,0.1292]$. Both retain their saved Holm-adjusted $p=0.0012$. These follow-ups improve fault-field precision conditional on their existing cohorts and masks, not generalization to new trained models.

The historical qPolarGrad code is preserved under its recovered definition. The extra factor in the old printed expression is documented rather than silently repaired. Its prior hardware-flow results are not reintroduced as matched evidence for the present checkpoints. Earlier FP32/BF16 studies, including the 96-bit language-model pilots and 1,024-bit DistilBERT classifier subset, remain separate restricted-universe evidence. They are not pooled with full-eligible INT8 curves or counted as independent architecture confirmations. Their representation effects and selection geometry require their own claims and controls.

\subsection{Remaining generalization requirements}
The main curves use one frozen checkpoint per model. Multiple fault fields quantify conditional fault randomness; they do not replace new checkpoints, independently sampled evaluation cohorts or repeated ranking construction. Published FKeras is missing from the transformer grids, and no matched PrioriFI-derived effectiveness grid is included. A stronger allocation claim requires a prespecified comparison against these feasible baselines, separate mask-construction and final-evaluation information, and uncertainty over the intended deployment populations. Those are prospective requirements, not unreported results.
